% Selected quotation and dialogue variants; fixed-font customization specimens.
\documentclass[aspectratio=169]{beamer}
\usetheme{bitpolito-studio}
\usepackage{bitpolitostudiotemplates}
\begin{document}
\studiotemplate[quote={A private matter is something one does not want the whole world to know, while a secret matter is something one does not want anyone to know. Privacy is the power to selectively reveal oneself to the world.},attribution={Eric Hughes, A Cypherpunk\textquotesingle{}s Manifesto, 1993}]{quotation}

\studiotemplate[variant=three-messages,message-one={Hurry up, you are going to be the next enlightened ruling class!},message-two={But you had never explained me how the society works and what are the values behind it...},message-three={It\textquotesingle{}s not needed.}]{dialogue}

\studiotemplate[variant=people,message={Let's talk about Bitcoin in our university!\\Of course only technical stuff...},caption={February 2020}]{dialogue}

\studiotemplate[variant=cards,message-one={Linux is a cancer},name-one={Ballmer},message-two={Open source is un-American},name-two={Allchin},message-three={Free software is communism},name-three={Gates},conclusion={GNU and Linux scared entrenched powers, but just\\worked better and developed a global community}]{dialogue}

\studiotemplate[quote={A shared system becomes easier to understand when participants can inspect its rules, compare different explanations and verify the evidence for themselves. Clear language helps people take part in that process.},attribution={A longer illustrative attribution, with context and a date}]{quotation}

\studiotemplate[variant=three-messages,message-one={A longer first message can explain\\why the conversation matters.},message-two={The reply can ask for more detail about the shared rules and the evidence available to everyone.},message-three={Let us verify it.}]{dialogue}

\studiotemplate[variant=people,message={Learning together starts with a question.\\Different views help everyone learn.},caption={An illustrative discussion with a longer context note}]{dialogue}

\studiotemplate[variant=cards,message-one={Discuss our\\assumptions},name-one={First participant},message-two={Compare different\\ways to participate},name-two={Second participant},message-three={Explain the next\\practical step},name-three={Third participant},conclusion={A longer conclusion connects the discussion to a shared next step that everyone can verify.}]{dialogue}

\studiotemplate[quote={},attribution={The attribution stays in its own field.}]{quotation}

\studiotemplate[quote={One idea can stand on its own.},attribution={}]{quotation}

\studiotemplate[variant=three-messages,message-one={},message-two={Only this response is visible.},message-three={}]{dialogue}

\studiotemplate[variant=people,message={},caption={The caption remains independent.},icon-person-two={}]{dialogue}

\studiotemplate[variant=cards,message-one={First point},name-one={},message-two={},name-two={Person 2},message-three={Third point},name-three={},conclusion={}]{dialogue}

\studiotemplate[variant=people,message={Each participant can use a different graphic.},caption={All replacements preserve their proportions.},icon-person-one={alice},image-person-two={assets/generated/wallet.pdf},graphic-person-three={\rule{3cm}{5mm}}]{dialogue}

\studiotemplate[variant=cards,message-one={Learn},name-one={Alice},icon-person-one={alice},message-two={Build},name-two={Bob},icon-person-two={bob},message-three={Verify},name-three={A tall graphic},graphic-person-three={\rule{5mm}{3cm}},conclusion={The bubbles and the concluding panel remain separate from the replacement graphics.}]{dialogue}

\studiotemplate[variant=cards,message-one={Learn},name-one={Person 1},message-two={Build},name-two={Person 2},message-three={Share},name-three={Person 3},conclusion={Hidden conclusion},show-conclusion=false]{dialogue}
\studiotemplate[quote={One clear idea.},attribution={Illustrative author}]{quotation}
\studiotemplate[quote={A short quotation.\\A deliberate second line.},attribution={A longer illustrative attribution that can wrap automatically to a second line while keeping its own fixed font size and the panel padding intact.}]{quotation}
\studiotemplate[quote={First line of a longer quotation.\\Second line of a longer quotation.\\Third line of a longer quotation.\\Fourth line of a longer quotation.\\Fifth line of a longer quotation.\\Sixth line of a longer quotation.\\Seventh line of a longer quotation.},attribution={An illustrative source, not a historical quotation}]{quotation}
\studiotemplate[quote={},attribution={}]{quotation}
\studiotemplate[width=compact,quote={One clear idea.},attribution={Illustrative author}]{quotation}
\studiotemplate[width=compact,quote={A shared system becomes easier to understand when participants can inspect its rules, compare explanations and verify the available evidence.},attribution={A longer illustrative attribution with context and a date}]{quotation}
\studiotemplate[width=medium,quote={One clear idea.},attribution={Illustrative author}]{quotation}
\studiotemplate[width=medium,quote={A shared system becomes easier to understand when participants can inspect its rules, compare explanations and verify the available evidence.},attribution={A longer illustrative attribution with context and a date}]{quotation}
\studiotemplate[width=wide,quote={One clear idea.},attribution={Illustrative author}]{quotation}
\studiotemplate[width=wide,quote={A shared system becomes easier to understand when participants can inspect its rules, compare explanations and verify the available evidence.},attribution={A longer illustrative attribution with context and a date}]{quotation}
\end{document}
